\section{附录 C：定理 \ref{thm:closure-dim} 的证明与二次域闭包限制}
\label{app:closure-dim-proof}

\subsection*{C.1 关系格与秩}

定义
\[
L'=\{(\mathbf{m},n)\in\ZZ^d\times\ZZ:\ \mathbf{m}\cdot\boldsymbol{\alpha}+n=0\}.
\]
向量 $(\alpha_1,\dots,\alpha_d,1)$ 在 $\QQ$ 上张成的维数为 $r(\boldsymbol{\alpha})$，故其整数关系格 $L'$ 的秩为 $(d+1)-r(\boldsymbol{\alpha})$。

\subsection*{C.2 投影同构}

投影映射 $(\mathbf{m},n)\mapsto \mathbf{m}$ 将 $L'$ 同构到
\[
L(\boldsymbol{\alpha})=\{\mathbf{m}\in\ZZ^d:\ \mathbf{m}\cdot\boldsymbol{\alpha}\in\ZZ\},
\]
因为 $\mathbf{m}\cdot\boldsymbol{\alpha}\in\ZZ$ 当且仅当存在 $n\in\ZZ$ 使 $\mathbf{m}\cdot\boldsymbol{\alpha}+n=0$。因此
\[
\mathrm{rank}\,L(\boldsymbol{\alpha})=\mathrm{rank}\,L'=(d+1)-r(\boldsymbol{\alpha}).
\]

\subsection*{C.3 闭包维数}

轨道闭包子环面维数满足
\[
\dim H=d-\mathrm{rank}\,L(\boldsymbol{\alpha})
=d-\bigl((d+1)-r(\boldsymbol{\alpha})\bigr)=r(\boldsymbol{\alpha})-1.
\]
从而若 $\alpha_i\in K$ 且 $[K:\QQ]=k$，则 $r(\boldsymbol{\alpha})\le k$，闭包维数 $\le k-1$。二次域 $k=2$ 时闭包维数 $\le 1$。

